\section{Introduction}

A global regularity certificate for a specified Navier--Stokes datum and a convergence theorem for a prescribed iteration answer different questions. The first controls the exact solution. The second must also control the iterates generated by a fixed starting rule. This paper establishes both kinds of result for explicit neighbourhoods of a nonlinear three-dimensional periodic field, while keeping their respective hypotheses separate.

We work with the unforced incompressible Navier--Stokes equations on
\(\mathbb T^3=(\mathbb R/2\pi\mathbb Z)^3\), with normalized spatial measure, viscosity one, and zero spatial mean. The centre is
\[
 W(x,y,z)=(\sin z,\sin x,\sin y).
\]
Its active wave vectors span three dimensions, and its projected quadratic nonlinearity is nonzero. In particular, its evolution is not simply its heat evolution. For a mean-zero field we use the Fourier norms
\[
 A_s(f)=\sum_{k\in\mathbb Z^3\setminus\{0\}}
       |k|^s|\widehat f(k)|_2,
 \qquad
 \|f\|_{\mathcal X}
   =\left(\int_0^\infty A_0(f(t))^2\,dt\right)^{1/2}.
\]
Here \(|\cdot|_2\) is the Hermitian Euclidean norm, also written \(|\cdot|\) below. We call \(\int_0^\infty A_0(u(t))^2\,dt\) the solution action. Unless stated otherwise, perturbations are smooth, real, mean-zero and divergence-free. They may have infinitely many Fourier modes. Smallness in \(A_{-1}\) does not impose a uniform bound on their higher Sobolev norms.

Our main result concerns the ordinary heat-seeded Picard sequence for initial data \(W+h\). With \(S(t)=e^{t\Delta}\) and the usual Leray-projected Duhamel operator \(Q\), this sequence is
\[
 v_h^{(0)}=S(W+h),\qquad
 v_h^{(n+1)}=S(W+h)-Q(v_h^{(n)},v_h^{(n)}).
\]
For \(A_{-1}(h)\le10^{-4}\), it converges in \(\mathcal X\) to the global smooth solution \(u_h\), and
\[
 \|v_h^{(n)}-u_h\|_{\mathcal X}
       \le 6(9/10)^{n-3},\qquad n\ge3.
\]
The proof certifies entry of the actual third iterate into a common invariant set. A finite nonnegative lower-triangular matrix then controls subsequent errors on a partition of the entire time half-line. The same construction gives quantitative dependence on the initial datum and a uniform Galerkin continuation bound. Thus the asserted rate concerns the specified ordinary iteration, rather than an iteration initialized at an independently chosen approximation.

A second result gives larger regularity neighbourhoods for data \(aW+h\). Among the certified regions are
\[
 \begin{split}
  |a|\le1,&\qquad A_{-1}(h)\le1/10,\\
  |a|\le11/10,&\qquad A_{-1}(h)\le1/40,\\
  |a|\le8/7,&\qquad A_{-1}(h)\le1/100.
 \end{split}
\]
These are separate regions, not the rectangle obtained by combining their largest amplitude and largest radius. They yield global smoothness and explicit continuation bounds, but the Picard convergence assertion above is not extended to their union.

Finally, a validated spherical Galerkin certificate at cutoff three provides a complementary estimate for the unperturbed datum. If \(u_M\) denotes an exact spherical Fourier--Galerkin solution starting from the projection of \(W\), the certificate gives
\[
 \sup_M\sup_{t\ge0}\|\nabla u_M(t)\|_2^2<14,
 \qquad
 \sup_M\int_0^\infty\|\nabla u_M(t)\|_2^4\,dt<21/2.
\]
The finite-dimensional residual includes every exterior mode generated by the reference trajectory. An analytic tail completes the estimate after the validated finite interval. These particular numerical bounds concern \(W\); they are not assigned unchanged to perturbed data.

The proof consists of analytic estimates together with finite arithmetic certificates. The certificates verify specific inequalities needed by the analytic argument; they do not replace the argument connecting Galerkin estimates to the full PDE. Our scope is the stated families and conclusions. We make no claim concerning unrestricted smooth initial data, optimal perturbation radii, or the earliest occurrence of the general methods used here.

For perspective, the zero-reference specialization of Lemma~\ref{ext:historylemma} gives global smoothness for smooth, real, mean-zero, solenoidal data with \(A_{-1}(u_0)<2\); its Galerkin-to-PDE justification is given after that lemma. Separately, Lemma~\ref{main:heat-lemma} gives the elementary Picard contraction condition \(A_{-1}(u_0)<1\), since \(4c\|Su_0\|_{\mathcal X}\le A_{-1}(u_0)\), with \(c=1/(2\sqrt2)\). The closed ball centred at zero with radius \(2\|Su_0\|_{\mathcal X}\) is invariant and strictly contractive and contains the heat seed. These are sufficient conditions, not optimal thresholds. Since \(A_{-1}(W)=3\), the zero-reference criterion covers \(aW+h\) whenever \(3|a|+A_{-1}(h)<2\), and an entire closed amplitude--radius class whenever \(3a_*+r<2\). By contrast, every datum of the main theorem satisfies \(A_{-1}(W+h)\ge3-10^{-4}>2\).

The reported action, stability and Galerkin bounds are certified estimates; the perturbation radii are sufficient admissible radii. Neither the bounds nor the radii are claimed optimal or representative of typical numerical behaviour.

\subsection{Relation to existing results}

Lei and Lin \cite{LeiLin2011} established global small-data mild solutions in a critical Fourier space on \(\mathbb R^3\). Ambrose, Lopes Filho and Nussenzveig Lopes \cite{Ambrose2024} prove periodic small-data well-posedness using a whole-half-line contraction in Fourier--Wiener spaces. Their argument also gives convergence from the ordinary heat seed in their space, hence in \(\mathcal X\): this last implication follows from
\[
 \|f\|_{\mathcal X}^2
 \le\big(\sup_{t\ge0}A_{-1}(f(t))\big)
       \int_0^\infty A_1(f(t))\,dt.
\]
Thus ordinary heat-seeded iteration and whole-half-line convergence are established principles. Here we supply the explicit neighbourhood of \(W\), certified entry of the actual third iterate, and the stated uniform bounds and rate.

Critical Fourier perturbation theory is also established. Benameur's Theorem~1.4 in \cite{Benameur2013} gives quantitative stability in \(\mathcal X^{-1}(\mathbb R^3)\) around an already global solution, with a perturbation threshold depending exponentially on a time-integrated Fourier norm of that solution. Its stated domain is \(\mathbb R^3\). Our periodic certificate supplies a specified reference, explicit verified bounds, and a convergence conclusion for a prescribed iteration. We do not present stability of global solutions as a new principle. Explicit high-Sobolev stability estimates for smooth periodic solutions are also part of the literature; see \cite{Pizzocchero2021}.

The use of approximate trajectories and invariant error tubes has precedent in the literature. Morosi and Pizzocchero \cite{Morosi2007}, Proposition~3.7 and Lemmas~4.8--4.9, construct solutions by Peano--Picard iteration from a suitable approximate seed. Their existence conclusion holds on the whole interval on which the control inequality is satisfied and can therefore be global; their stated iteration estimates give uniform convergence on each compact subinterval. Their later Sobolev control theorem \cite{Morosi2011}, Proposition~4.5, bounds an exact periodic Euler or Navier--Stokes solution through a scalar control inequality. These results provide the comparison framework. The task here is the explicit verification of the stated data classes and of the heat-seed entry and whole-half-line convergence estimates, not the introduction of a posteriori analysis itself.

Morosi, Pernici and Pizzocchero \cite{Morosi2014} use high-order Reynolds expansions and a posteriori Sobolev control for structured periodic initial data. Their Proposition~2.3 gives the analytic implication from a control equation to exact-solution bounds. The paper explicitly distinguishes its numerical control-equation calculations from certified computer-assisted proofs in footnote~4. A Reynolds expansion is not the ordinary heat-seeded Picard sequence. Accordingly, the relevant comparisons concern the approximation rule, function space, data class and certification standard, rather than only a displayed amplitude or Reynolds number.

Silva \cite{Silva2026}, Theorem~2.2, states a global smoothness theorem on the same normalized torus at viscosity one for the cyclic-shear class
\[
 \mathcal C_S=
 \{aU+h:0\le a\le7/10,\ A_2(h)\le1/100\},
 \qquad U=(\sin y,\sin z,\sin x),
\]
where the perturbations satisfy the same smoothness, reality, mean and divergence conditions. The theorem also supplies a residual-action estimate uniform over cubical Galerkin cutoffs. That action estimate is a distinct conclusion from the bounds stated here.

For clarity, the inclusion between the corresponding regularity classes can be checked directly. Let \(P\) be the orthogonal matrix interchanging the first two coordinates and define
\[
 (\mathcal T f)(x)=P f(P^T x).
\]
This simultaneous coordinate and velocity-component permutation preserves the Navier--Stokes equations and every \(A_s\) norm, and \(\mathcal T U=W\). Since \(|k|\ge1\) for every nonzero integer wave vector,
\[
 A_{-1}(\mathcal T h)\le A_2(\mathcal T h)=A_2(h).
\]
Consequently, \(\mathcal T\mathcal C_S\) is contained in our region
\[
 \mathcal C_E=
 \{aW+h:|a|\le8/7,\ A_{-1}(h)\le1/100\}.
\]
The inclusion is strict. Indeed, \(W\in\mathcal C_E\), whereas any representation \(W=bW+h\), with \(0\le b\le7/10\), necessarily has
\[
 A_2(h)=A_2((1-b)W)=3(1-b)\ge9/10>1/100.
\]
Thus \(W\notin\mathcal T\mathcal C_S\), even allowing an alternative amplitude decomposition.

This comparison establishes a strict enlargement of this specified cyclic-shear class for the conclusion of global PDE regularity. It neither compares the numerical residual-action constants nor transfers estimates between cubical and spherical cutoffs. It also does not assert inclusion of Silva's separate ABC, TG3 or coefficient families. The main Picard theorem supplies a different additional conclusion on its own smaller neighbourhood. Attribution and development history are separate from these mathematical inclusion statements.
